\documentclass[aspectratio=169]{beamer}
\usepackage{array}
%\usepackage{dirtytalk}
\usepackage[dvipsnames]{xcolor}
\title{IANA YANG~SIDs Registry, Constrained~YANG~Library, Unified~Datastore}
\subtitle{CoRE Interim 07}
\author{Vojtěch Vilímek}
\date{2026-05-20}

% source: https://tex.stackexchange.com/questions/137022/how-to-insert-page-number-in-beamer-navigation-symbols
\setbeamertemplate{sidebar right}{}
\setbeamertemplate{footline}{%
\hfill\usebeamertemplate***{navigation symbols}
\hspace{1cm}\insertframenumber{}/\inserttotalframenumber}

\hypersetup {
  colorlinks=true,
  urlcolor=MidnightBlue,
}

\begin{document}

\begin{frame}
  \titlepage
\end{frame}

\begin{frame}
  \frametitle{Agenda}

  \begin{itemize}
    \item IANA YANG SIDs Registry
    \item Constrained YANG Library
    \item Unified Datastore Semantics
    \item PYCORECONF Discussion
  \end{itemize}	
\end{frame}

\section{IANA YANG SIDs Registry}
\begin{frame}
  \frametitle{IANA YANG SIDs Registry} 

  \begin{block}{\href{https://mailarchive.ietf.org/arch/msg/yang-tooling/Sr-AMBFO-WBoHcIog7iaMDqS2I4/}{Andy Bierman's comments}}
    \begin{itemize}
      \item no YANG Doctors review established for .sid Files (or YD experts review).
      \item iana-if-types: allocation too small
      \vspace{0.3em}

      SID range has 400 numbers, allocation for revision 2026-03-17 leaving 94 SIDs unassigned.

      Andy would like to have 10000 SID range.

      Could be solved by additional SID range when necessary.

      \item ietf-system, \href{https://datatracker.ietf.org/doc/html/rfc9595}{RFC9595} example is not correct
      RFC9595 .sid file example should not be considered normative.

      \item ietf-yang-types and ietf-inet-types: old reference to 
      \href{https://datatracker.ietf.org/doc/html/rfc6991}{RFC6991} not \href{https://datatracker.ietf.org/doc/html/rfc9911}{RFC9911} -- Fixed

      \item (NIT) ietf-coreconf entry has reference to RFC9595.
    \end{itemize}
  \end{block}
\end{frame}

\begin{frame}
  \frametitle{IANA YANG SIDs Registry}
  \href{https://www.iana.org/assignments/yang-sid/yang-sid.xhtml}{Link.}

  Errata wording then Registry population.

  General: Should we use size 50 for typedef only modules?

  \vspace{0.3em}
  \begin{quote}
    File ietf-yang-types@2025-12-22.sid created

    Number of SIDs available : 50

      Number of SIDs used : 1
  \end{quote}

  \vspace{1em}
  \begin{quote}
    File ietf-inet-types@2025-12-22.sid created

    Number of SIDs available : 40

    Number of SIDs used : 1
  \end{quote}
\end{frame}

\begin{frame}
  \frametitle{IANA YANG SIDs Registry}

  Should we consider automatic output as good enough for standard .sid files?

  \vspace{1em}
  Current modules in registry: 

    iana-crypt-hash, iana-if-types, ietf-coreconf, ietf-ient-types, ietf-interfaces,
    ietf-ip, ietf-netconf-acm, ietf-schc, ietf-sid-file, ietf-system,
    ietf-voucher-request, ietf-voucher, ietf-yang-types

  \vspace{1em} 
  Why ietf-sid-file? Only standardized encoding is JSON.
\end{frame}

\section{Constrained YANG Library}
\begin{frame}[fragile]
  \frametitle{Constrained YANG Library}

\begin{verbatim}
// file ietf-constrained-yang-library.yang
  list import-only-module {
    key "identifier revision";
    leaf revision {
      type union {
      type revision-identifier;
        type string; // would not be empty better?
          legnth 0;
        }
      }
      description
        "The YANG module revision date.";
    }
  }
\end{verbatim}

  Need to allocate SID range, and create .sid File.
\end{frame}


\section{Unified DS Semantics}
\begin{frame}[fragile]
  \frametitle{Unified Datastore Semantics}
  \begin{block}{2.4.  Unified datastore}

CORECONF supports a simple datastore model consisting of 
a single unified datastore.  This datastore provides 
access to both configuration and operational data.  
Configuration updates performed on this datastore are 
reflected immediately or with a minimal delay as 
operational data.

More complex datastore models such as the Network 
Management Datastore Architecture (NMDA) as defined 
by [\href{https://datatracker.ietf.org/doc/html/rfc8342}{RFC8342}] are out of scope of the present
specification.
  \end{block}

\end{frame}

\begin{frame}
  \frametitle{Unified Datastore Semantics}
Characteristics of the unified datastore are summarized in the table
below:

  \vspace{1em}
  \begin{tabular}{ | l | l | }
    \hline
    \textbf{Name}         &  \textbf{Value} \\ 
    \hline
    Name         &  unified \\
    \hline
    YANG modules &  all modules \\
    \hline
    YANG nodes   & all data nodes ("config true" and "config false") \\
    \hline
    Access       & read-write \\
    \hline
    How applied  &  changes applied in place immediately or with a    
                    minimal delay \\
    \hline
    Protocols    & CORECONF \\
    \hline
    Defined in   & "ietf-coreconf" \\
    \hline
  \end{tabular}
\end{frame}

\begin{frame}
  \frametitle{Unified Datastore Semantics}

  \newlength{\fNameWidth}
  \settowidth{\fNameWidth}{YANG modules}
  \newlength{\fValueWidth}
  \settowidth{\fValueWidth}{chagnes applied in place immediately or with minimal delay}
  \begin{tabular}{ | m{\fNameWidth} | m{\fValueWidth} | }
    \hline
    Access & read-write \\
    \hline
  \end{tabular}
  \vspace{0.5em}

  Are all data nodes read-write, or just the "config true"? The "config false" are read-only.

  \vspace{1em}
  \begin{tabular}{ | m{\fNameWidth} | m{\fValueWidth} | }
    \hline
    How applied & changes applied in place immediately or with a
                  minimal delay \\
    \hline
  \end{tabular}
  \vspace{0.5em}

  Do I have two value, where the "config true" shadows the "config false", or just one value?  
\end{frame}


\begin{frame}
  \frametitle{Unified Datastore Semantics}

  \begin{block}{Consistency}
    Operational data are always changing.

    For more complex queries, data consistency may be broken during message processing.

    Also for iPATCH requests, rollback functionality is requered for two-or-more
    elements of CBOR sequence.

    \vspace{1em}
    Isn't the rollback too resource heavy for the constrained environment?
  \end{block}
\end{frame}

\begin{frame}[fragile]
  \frametitle{CORECONF Consistency}

  \begin{block}{Examples}
\footnotesize \begin{verbatim}
REQ: FETCH </c>  (Content-Format: application/yang-identifiers+cbor-seq)
[1533, "eth0"]   / interface (SID 1533) with name = "eth0" /
1723,            / current-datetime (SID 1723) /

RES: 2.05 Content  (Content-Format: application/yang-instances+cbor-seq)
{
  1533 : {
     4 : "eth0",              / name (SID 1537) /
    13 : "referenced-value"    / leafref to 1723 /
  }
}
/ in the meantime the instance 1723 caesed to exist /
{
  1723 : null          / current-datetime (SID 1723) /
},
\end{verbatim}
  \end{block}
\end{frame}

\begin{frame}[fragile]
  \frametitle{CORECONF Consistency}

  \begin{block}{Examples}
\footnotesize \begin{verbatim}
REQ: iPATCH </c>   (Content-Format: application/yang-instances+cbor-seq)
{  1755 : true  }                / enabled (SID 1755) /
[  "obvious", "error" ]
{  1756 : {                      / server (SID 1756)  /
     3 : "tic.nrc.ca",           / name (SID 1759)    /
     4 : true,                   / prefer (SID 1760)  /
     5 : {                       / udp (SID 1761)     /
       1 : "132.246.11.231" }    / address (SID 1762) /
}  }

RES: 4.00 Bad Request
{ 1025: {
    4: 1019  / error-tag     (SID 1029) -> operation-failed  (SID 1019) /
    1: 1012  / error-app-tag (SID 1026) -> malformed-message (SID 1012) /
    3: "Expected map, got array"    / error-message (SID 1028) /
  }
}
\end{verbatim}
  \end{block}
\end{frame}


\section{PYCORECONF Discussion}
\begin{frame}
  \begin{block}{PYCORECONF Discussion}
    \begin{itemize}
      \item YANG-oblivious process
      \item all built-in types
      \item typedefs
      \item maybe work on sid-file-ext.yang?
    \end{itemize}
  \end{block}
\end{frame}

\begin{frame}
  \begin{block}{PYCORECONF Discussion}
  \end{block}
\end{frame}

\end{document}
